\metatitle{Schur P- and Q-functions}
\metadescription{Schur P- and Q-functions, shifted tableaux, Pfaffian and
Cauchy identities, and generalized Schur P- and Q-functions.}
\metakeywords{Schur P function, Schur Q function, shifted tableau, Pfaffian,
generalized Schur P function}

\section[schurP]{Schur's $P$ functions}


\begin{polydata}{schurP}
  Name   & Schur's $P$ functions \\
  Space    & Sym \\
  Basis    & False \\
  Rating   & 4 \\
  Bib      & Schur1911\\
  Year     & 1911 \\
  Keywords & pfaffian,tableaux,schur-positive \\
  Symbol   & $\schurP_\lambda(\xvec)$ \\
  Category & Schur \\
  PositiveIn & schurS | Stembridge1989 \\
  TransformsTo & schurQ | Schur1911 | map=multiply by $2^{\length(\lambda)}$ \\
\end{polydata}

Schur's $P$ functions were introduced by \name{Issai Schur} in \cite{Schur1911}
to study projective representations of the symmetric and alternating groups.
Schur's $P$ functions $\schurP_\lambda(\xvec)$ are indexed by
\defin{strict partitions}\label{strictPartition}, that is, partitions with
distinct parts, and are given
as the specialization of the
\hyperref[hallLittlewoodP]{Hall--Littlewood $P$ functions} at $t=-1$.
\name{Federico Ardila} and \name{Luis G. Serrano} prove Stanley's conjecture
that, for the staircase shape $\delta_n$, every skew Schur function
$\schurS_{\delta_n/\mu}$ expands nonnegatively in Schur's $P$ functions
\cite{ArdilaSerrano2012}.  Their coefficients count fillings of shifted
shapes; the special case $\schurS_{\delta_n/\delta_{n-2}}$ is related to
\hyperref[eulerianSymmetry]{Eulerian numbers and alternating permutations}.
\name{Matthew Fayers} proves a Pieri-type rule for Schur's $P$ functions in
\cite[Prop.~3.6]{Fayers2020}.  For the reduction operators
$\partial_{(a)}$ and $\partial_{(1^a)}$, one has
\[
\partial_{(a)}\schurP_\lambda=\partial_{(1^a)}\schurP_\lambda
=\sum_\mu Hs_{\lambda/\mu}\schurP_\mu,
\]
where $Hs_{\lambda/\mu}=2^N$ if $\lambda/\mu$ is a horizontal strip and $N$
counts columns $c+1$ containing a box of $\lambda/\mu$ such that column $c$
contains none; otherwise $Hs_{\lambda/\mu}=0$.


\subsection[schursPPfaffian]{Pfaffian formula}

Let $\lambda$ be a partition with $\ell$ distinct parts and let $n \geq \ell$.
Then
\[
\schurP_\lambda(x_1,\dotsc,x_n)
\coloneqq
\frac{1}{(n-\ell)!} \sum_{\sigma \in \symS_n} \sigma \left(
\xvec^{\lambda} \prod_{\substack{i \leq \ell \\  i \lt j \leq n } } \frac{x_i+x_j}{x_i-x_j}
\right).
\]

\subsection[schursPTableauformula]{Combinatorial formula}

Given a partition $\lambda$ with distinct parts, consider the
\defin{shifted tableau}\label{shiftedTableau} of shape $\lambda$.
We consider fillings of this diagram with entries in the alphabet
\[
1' \lt 1 \lt 2' \lt 2 \lt \dotsb
\]
such that
\begin{itemize}
\item
	each row has at most one marked $i$ for every $i=1,2,\dotsc$
\item
	each column has at most one unmarked $i$, for every $i=1,2,\dotsc$,
\item
	entries in rows and columns are weakly increasing and
\item
	there are no marked entries on the main diagonal.
\end{itemize}
We let $ShSSYT(\lambda)$ denote the set of such fillings.

\begin{example}
For $\lambda=9531$ we have the shifted diagram with an element of
$ShSSYT(\lambda)$.
\begin{figure}
\begin{ytableau}
 &   &  & & & & & & \\
\none & & & & & \\
\none & \none & & & \\
\none & \none & \none &
\end{ytableau}
\begin{ytableau}
1 & 1  & 2' & 3' & 3 & 3 & 4 & 4 & 4 \\
\none & 2 & 2 & 3 & 4' & 4 \\
\none & \none & 2 & 3 & 4 \\
\none & \none & \none & 4
\end{ytableau}
\end{figure}
\end{example}

\defin{Schur's $P$ function}\label{schurPFunction} is defined by
\[
\schurP_\lambda(\xvec) \coloneqq \sum_{T \in ShSSYT(\lambda)} \xvec^T
\]
where $\xvec^T$ is the product over all entries in $T$, and $i$ and $i'$
both contribute $x_i$.

\name{Ronald C. King} and \name{A. M. Hamel} study Hall--Littlewood
polynomials $P_\lambda(x;t)$ attached to arbitrary simple Lie algebras at
$t=-1$ \cite{KingHamel2007}.  In type $A$, the case of partitions with
distinct parts recovers Schur's $P$ functions.  For type $C$, they give
analogous $Q$-functions using primed symplectic shifted tableaux, together
with Weyl-group symmetry and lattice-path determinant interpretations.

An alternative combinatorial model uses semistandard decomposition tableaux,
introduced in \cite{Serrano2010}.
See \cite{Cho2012} for additional results.

\subsection[schurPSchurExpansion]{Schur expansion}

Schur's $P$ functions are \hyperref[schurPositiveDefinition]{Schur-positive};
see \cite{WorleyThesis,Sagan1987,Assaf2018ShiftedDual}.
\name{Bruce Sagan} gives a proof using an insertion algorithm similar to
\hyperref[rsk]{RSK} insertion that respects the
\hyperref[knuthRelations]{Knuth relations}.
The proof by \name{Sami Assaf} uses
\hyperref[schurPositivityDualEquivalence]{dual equivalence}.

\name{Seung-Il Choi} and \name{Jae-Hoon Kwon} give a proof of
\hyperref[schurPositiveDefinition]{Schur positivity} using
\hyperref[crystals]{crystals}; see \cite{ChoiKwon2018}.
A different proof using crystals is provided in \cite{AssafOguz2018}.


\begin{example*}[Schur expansions]
Here are the Schur expansions of $\schurP_\lambda(\xvec)$ for various $\lambda$.

\begin{array}{lll}
\toprule
\text{Partition} & \text{Schur expansion} \\
\midrule
1	& \schurS_{1} \\
2	& \schurS_{2} + \schurS_{11} \\
21	& \schurS_{21} \\
3	& \schurS_{3} + \schurS_{21} + \schurS_{111} \\
31	& \schurS_{22} + \schurS_{31} + \schurS_{211} \\
321	& \schurS_{321} \\
4	& \schurS_{4} + \schurS_{31} + \schurS_{211} + \schurS_{1111} \\
41	& \schurS_{32} + \schurS_{41} + \schurS_{221} + \schurS_{311} + \schurS_{2111} \\
32	& \schurS_{32} + \schurS_{221} + \schurS_{311} \\
421	& \schurS_{322} + \schurS_{331} + \schurS_{421} + \schurS_{3211} \\
\bottomrule
\end{array}
\end{example*}

\subsection[schurPLRRule]{Littlewood--Richardson rule}

The coefficients $f_{\lambda\mu}^{\nu}$ in
\[
\schurP_\lambda(\xvec) \schurP_\mu(\xvec) = \sum_{\nu} f_{\lambda\mu}^{\nu} \schurP_\nu(\xvec)
\]
are nonnegative integers; see \cite{Stembridge1989}.  A combinatorial
description of these coefficients
was first given by \name{Luis Serrano} \cite{Serrano2009}.
Another proof is given in \cite{Cho2012}.
\name{Sami Assaf} gave a proof using
\hyperref[schurPositivityDualEquivalence]{dual equivalence}
\cite{Assaf2018ShiftedDual}.

The coefficients $f_{\lambda\mu}^{\nu}$ appear in the study of the type $B$
and type $D$ \hyperref[schubertCalculus]{Schubert calculus} of
\hyperref[grassmannianVarieties]{Grassmannians}.

\name{Duc Khanh Nguyen} gives a new combinatorial model for
$f_{\lambda\mu}^{\nu}$ and for the coefficients in the Schur expansion of
$\schurQ_\lambda$ \cite{Nguyen2022}.
\name[Santiago Estupinan-Salamanca]{Santiago Estupiñán-Salamanca} and
\name{Oliver Pechenik} give an easier formula for $f_{\lambda\mu}^{\nu}$
using a type of Yamanouchi word \cite{EstupinanSalamancaPechenik2025x}.


\subsection[schurPEvacuation]{Evacuation}

For shifted tableau switching and related \hyperref[evacuation]{evacuation} operations, see
\cite{ChoiNamOh2017x}.


\section[schurQ]{Schur's $Q$ functions}


\begin{polydata}{schurQ}
  Name   & Schur's $Q$ functions \\
  Space    & Sym \\
  Basis    & False \\
  Rating   & 3 \\
  Bib      & Schur1911\\
  Year     & 1911 \\
  Keywords & pfaffian,tableaux,schur-positive \\
  Symbol   & $\schurQ_\lambda(\xvec)$ \\
  Category & Schur \\
  PositiveIn & schurS | Stembridge1989 \\
\end{polydata}


\defin{Schur's $Q$ functions}\label{schurQFunction}
$\schurQ_\lambda(\xvec)$ are indexed by strict partitions
and are given
as the specialization of the
\hyperref[hallLittlewoodQ]{Hall--Littlewood $Q$ functions} at $t=-1$.
\name{Richard Stanley} gives a combinatorial formula using
\hyperref[gtpatterns]{Gelfand--Tsetlin patterns} \cite{Stanley1986}.

Equivalently,
$\schurQ_\lambda(\xvec) = 2^{\length(\lambda)} \schurP_\lambda(\xvec)$.
Thus $\schurQ_\lambda(\xvec)$ can be realized as a sum over tableau objects
by modifying the formula for Schur's $P$ functions and allowing diagonal
entries in $ShSSYT(\lambda)$ to be marked.


Schur's $Q$ functions are related to
\defin{spin characters}\label{spinCharacter},
which were part of the original motivation for studying these functions.
The following is due to Schur.

\begin{theorem}[See \cite[3.3.6]{Morris1962} and \cite[(7.1)]{Stembridge1989}]

Let $\lambda$ be a partition of $n$ with $m$ distinct parts.
The spin characters $\zeta_{\mu}^{\lambda}$ satisfy
\[
\schurQ_\lambda(\xvec) = \sum_{\mu}
2^{(\length(\lambda)+ \length(\mu)+\epsilon)/2}
\zeta_{\mu}^{\lambda} \frac{ \powerSum_{\mu} }{z_\mu} .
\]
Here, $\mu$ runs over partitions with odd parts,
and $\epsilon$ is $0$ or $1$ depending on whether $n-m$ is even or odd.
\end{theorem}

The one-part partitions have a rather explicit expression:
\[
\schurQ_{(n)}(\xvec) = \sum_{\lambda \vdash n} 2^{\length(\lambda)} \monomial_{\lambda}(\xvec).
\]


\name{Hjalmar Rosengren} gives formulas for
$\schurQ_{\lambda}(1,q,q^2,\dotsc,q^n)$ and also shows
\cite[Prop.~3.1]{Rosengren2008} that for a partition $\lambda$ of length $m$,
\[
\schurQ_{\lambda}(1,q,q^2,\dotsc) = \prod_{i=1}^m \frac{(-1;q)_{\lambda_i}}{ (q;q)_{\lambda_i} }
\prod_{1 \leq i \lt j \leq m} \frac{q^{\lambda_j} - q^{\lambda_i}}{1-q^{\lambda_i+\lambda_j}}.
\]
Rosengren notes that this formula is equivalent to
\hyperref[schurKawanaka]{Kawanaka's identity} for Schur functions.


A \hyperref[murnaghanNakayamaRule]{Murnaghan--Nakayama}-type formula for computing
the spin characters $\zeta_{\mu}^{\lambda}$ is given in \cite{MorrisOlsson1988}.

\name{Alun Morris} and \name{Jørn Olsson} give \cite[Cor.~4.3]{MorrisOlsson1988}
a formula which seems to be a spin
version of the \hyperref[borderStripHookFormulas]{Fomin--Lulov}
hook formula for ribbon tableaux.
This is similar to the properties of the classical irreducible
character values for the symmetric group.
Morris and Olsson use the abacus model to prove their results.


A Hall--Littlewood analogue of $\schurQ_\lambda(\xvec)$ was introduced in
\cite{TudoseZabrocki2003}.

\name{Soojin Cho}, \name{JiSun Huh}, and \name{Sun-Young Nam} conjecture a
classification of ribbon shapes for which the functions
$\schurQ_{\lambda/\mu}(\xvec)$ are \hyperref[powerSum]{power-sum} positive
\cite{ChoHuhNam2020}.

\name{Yuta Nishiyama} studies a conjecture of
\name{Hiroyuki Mizukawa}, \name{Masatoshi Nakajima}, and
\name{Hiro-Fumi Yamada} relating 2-reduced Schur functions and
Schur $Q$-functions \cite{Nishiyama2022x}.  The conjecture expresses certain
sums involving Littlewood--Richardson coefficients and two 2-reduced Schur
functions as scalar multiples of Schur $Q$-functions.

For more on spin characters, see \cite{MatsumotoSniady2020}.

\hyperref[plethysm]{Plethysm} stability for Schur's $Q$-functions is studied
in \cite{GrafJing2024x}.

\name{John Graf} and \name{Naihuan Jing} give a Pfaffian formulation of
Schur's $Q$-functions which allows the indexing composition to have negative
parts \cite{GrafJing2025}.  This makes the Pfaffian construction compatible
with the tableau and vertex-operator pictures and gives algebraic proofs of
several standard $Q$-function identities.
\name{S.-J. Lee} studies shifted $t$-Schur functions arising from the
modified odd Greaves--Jing--Zhu operator
\cite{Lee2026x,Lee2026xModifiedOperator,Lee2026xTransitionMatrices}.  These functions
are obtained from Schur's $Q$-functions by the plethystic substitution
$\mathcal{Q}_\lambda(X;t)=Q_\lambda[X-tX]$, so the associated weight
$\mathcal{Q}_\lambda(X;t)P_\lambda(Y)$ is a shifted Schur weight with a
virtual first alphabet.  The paper gives the normalization, Pfaffian
correlation kernel, Fredholm Pfaffian for the largest part, and size
cumulants; at $t=-q$ with $q\geq 0$ the virtual alphabet becomes the positive
alphabet $X+qX$.
The companion papers establish two-row and Pfaffian Giambelli formulas,
Cauchy and finite shifted-Gessel identities, and explicit transition matrices
between the shifted $t$-Schur bases.  At cyclotomic specializations those
matrices have Schur-$Q$-positive entries, together with reciprocity, rank, and
unimodality properties.

\name{Colin Defant} and \name{Dominic Searles} introduce shifted domino
functions, a type-$B$ analogue of Schur's $Q$-functions, using type-$B$
$0$-Hecke modules and standard shifted domino tableaux
\cite{DefantSearles2024x}.  These functions expand positively in
\name{T. Kyle Petersen}'s type-$B$ peak functions.



\subsection[schurQCauchy]{Cauchy identity}

\name{Bruce Sagan} proves \cite[Cor.~8.3]{Sagan1987} that
\begin{equation*}
\prod_{i,j=1}^\infty \frac{1+x_iy_j}{1-x_iy_j}
=
\sum_{\lambda} 2^{-\length(\lambda)}
\schurQ_\lambda(\xvec)\schurQ_\lambda(\yvec),
\end{equation*}
where the sum is over all strict partitions.


\section[generalizedSchurP]{Generalized Schur $P$- and $Q$-functions}

\begin{polydata}{generalizedSchurP}
  Name     & Generalized Schur $P$-functions \\
  Space    & Sym \\
  Basis    & False \\
  Rating   & 1 \\
  Bib      & Okada2019x \\
  Year     & 2019 \\
  Keywords & pfaffian,ninth-variation \\
  Symbol   & $P^F_\lambda(\xvec)$ \\
  Category & Schur \\
  Generalizes & schurP | Okada2019x | scope=$f_d(u)=u^d$ specialization \\
\end{polydata}

\begin{polydata}{generalizedSchurQ}
  Name     & Generalized Schur $Q$-functions \\
  Space    & Sym \\
  Basis    & False \\
  Rating   & 1 \\
  Bib      & Okada2019x \\
  Year     & 2019 \\
  Keywords & pfaffian,ninth-variation \\
  Symbol   & $Q^F_\lambda(\xvec)$ \\
  Category & Schur \\
  Generalizes & schurQ | Okada2019x | scope=$f_d(u)=2u^d$ specialization \\
\end{polydata}

\name{Soichi Okada} introduces
\defin{generalized Schur $P$-functions}\label{generalizedSchurPFunction} and
\defin{generalized Schur $Q$-functions}\label{generalizedSchurQFunction}
associated with an admissible polynomial sequence $F=\{f_d\}_{d\geq 0}$
\cite{Okada2019x}.  The functions $P^F_\lambda(\xvec)$ are defined by a
Nimmo-type Pfaffian formula, and this construction is Macdonald's ninth
variation in the $P$/$Q$ setting.  The ordinary Schur $P$- and $Q$-functions
are recovered by taking $f_d(u)=u^d$ and $f_d(u)=2u^d$, respectively.
The same framework includes Ivanov's factorial $P$- and $Q$-functions and
the $t=-1$ Hall--Littlewood functions attached to classical root systems.
Okada proves Cauchy-type identities, a J{\'o}zefiak--Pragacz--Nimmo formula
for skew $Q$-functions, and a Pieri-type rule.
